\section{Simulation Study}\label{Sec: simulation}

We investigate three questions: whether the distributed estimator retains the
accuracy of pooled KuLSIF, how optimization accuracy depends on communication,
and how heterogeneous client distributions and unequal sample sizes affect
estimation. The experiments distinguish the statistical performance of the
pooled estimator from the effects of random-feature approximation and finite
optimization.

\paragraph{Data-generating model.}
We consider $C=4$ clients and $d=5$ covariates. Let
$F_c=\mathcal{N}_5(\mu_c,I_5)$, where
\[
\mu_1=(-1,-1,0,0,0)^\top,\quad
\mu_2=(-1,1,0,0,0)^\top,\quad
\mu_3=(1,-1,0,0,0)^\top,\quad
\mu_4=(1,1,0,0,0)^\top.
\]
Unless otherwise specified, client $c$ samples from $P_c=F_c$, the source
proportions are $\pi_c=1/4$, and
\[
P=\sum_{c=1}^4\pi_cF_c,\qquad
Q=\sum_{c=1}^4\rho_cF_c,\qquad
(\rho_1,\rho_2,\rho_3,\rho_4)=(0.05,0.15,0.30,0.50).
\]
Writing $f_c$ for the density of $F_c$, the true ratio is
\[
r^*(x)=\frac{\sum_{c=1}^4\rho_cf_c(x)}
                 {\sum_{c=1}^4\pi_cf_c(x)}.
\]
This construction yields a nonlinear density ratio that can be evaluated
exactly, allowing us to assess estimation accuracy against the known truth.

\paragraph{Methods and parameter selection.}
We compare distributed gradient descent (GD; Algorithm~\ref{alg:gd}),
closed-form random-feature KuLSIF (RF; equation~\eqref{eq: RF-linear-coef}),
and pooled Gaussian-kernel KuLSIF (Kernel). The pooled kernel method is an
accuracy reference requiring centralized source data. We use
$K_\sigma(x,x')=\exp\{-\|x-x'\|_2^2/(2\sigma^2)\}$ and its cosine random
features, with $R=200$ unless stated otherwise. GD starts at zero and uses
$\eta=(2+\lambda)^{-1}$, corresponding to the uniform bound of $2$ on the
squared cosine feature functions. Its server receives only client
matrix--vector products, not source observations or second-moment matrices.
All reported predictions are untruncated.

Let $N$ and $M$ denote the total available source and target sample sizes.
Each client's source sample and the target sample are split into $80\%$
fitting data and $20\%$ validation data. Thus the sample sizes entering the
estimators are $n=0.8N$ and $m=0.8M$, and the client weights are the fitting
proportions $\pi_c=n_c/n$. The default setting is $(N,M)=(2000,1000)$.
For the sample-size and unequal-size comparisons, each KuLSIF method selects
its parameters separately by minimizing
\[
\widehat{\mathcal V}(r)
=\frac12\sum_{c=1}^C\frac{\pi_c}{v_c}
  \sum_{i=1}^{v_c}r(X_{ci}^{\mathrm{val}})^2
-\frac{1}{v_Q}\sum_{j=1}^{v_Q}r(Z_j^{\mathrm{val}}),
\]
where $v_c$ and $v_Q$ are the corresponding validation sample sizes.
We search $\lambda\in\{10^{-4},10^{-3},10^{-2},10^{-1},1\}$ and
$\sigma\in\{0.5,1,2\}\sigma_0$, where $\sigma_0$ is the median pairwise
distance among the fitting target observations. GD uses $5000$ updates for
each candidate and $10000$ updates for the final fit at the selected
parameters, using only the fitting observations. RF and GD share the same
feature-generation randomness within each replication. The optimization and
controlled-heterogeneity experiments instead use matched parameters, as
specified below, to isolate their respective mechanisms.

\paragraph{Evaluation and Monte Carlo uncertainty.}
For each replication, we generate an independent stratified test sample
$\{X_{cs}^{\mathrm{test}}:s\in[S_c]\}$ from the client laws, with
$\sum_cS_c=20000$ and $S_c$ proportional to $\pi_c$. We report
\[
\operatorname{RMSE}(\widehat r)
=\left\{\sum_{c=1}^C\frac{\pi_c}{S_c}
\sum_{s=1}^{S_c}
\bigl[\widehat r(X_{cs}^{\mathrm{test}})
      -r^*(X_{cs}^{\mathrm{test}})\bigr]^2\right\}^{1/2},
\]
which approximates $\|\widehat r-r^*\|_{L^2(P)}$.
Methods share fitting, validation, and test samples within each replication.
The tables report the mean of the replicate RMSEs, with Monte Carlo standard
errors (MCSEs) in parentheses. We use $100$ replications for each controlled
mechanism setting and $200$--$500$ for each tuned setting. Counts for the
latter were determined from separate pilot variance estimates, targeting a
paired MCSE of $0.0025$ for GD versus the RF and Kernel references, with a
cap of $500$; the cap does not guarantee this precision in every setting.
The $240$ pilot datasets are excluded from the $3306$ production datasets.
No numerical fitting failures occurred. On three independent integration-check
datasets, increasing the test size from $20000$ to $100000$ changed RMSE
by less than $0.0006$ for the estimators checked.

\subsection{Estimation Accuracy as Sample Sizes Increase}
\label{Sec: simulation-accuracy}

We vary $N\in\{1000,2000,4000,8000\}$ with $M=1000$ and
$M\in\{250,500,1000,2000\}$ with $N=2000$. The default setting is shared
between these comparisons. The full-kernel and random-feature references
are evaluated at every sample size.

\begin{table}[htbp]
\centering
\caption{Estimation accuracy: mean independent-test RMSE with MCSE in
parentheses. $N$ and $M$ include validation observations; $L$ is the number
of production replications.}
\label{tab:simulation-accuracy}
\small
\begin{tabular}{rrrccc}
\toprule
$N$ & $M$ & $L$ & GD & RF & Kernel \\
\midrule
1000 & 1000 & 500 & 0.2328 (0.0036) & 0.2292 (0.0031) & 0.2273 (0.0033) \\
2000 & 1000 & 306 & 0.2075 (0.0024) & 0.2064 (0.0026) & 0.2061 (0.0028) \\
4000 & 1000 & 200 & 0.2032 (0.0032) & 0.2020 (0.0033) & 0.2023 (0.0035) \\
8000 & 1000 & 200 & 0.1947 (0.0027) & 0.1939 (0.0027) & 0.1906 (0.0026) \\
\midrule
2000 & 250  & 500 & 0.2836 (0.0047) & 0.2790 (0.0044) & 0.2791 (0.0045) \\
2000 & 500  & 500 & 0.2356 (0.0029) & 0.2326 (0.0028) & 0.2313 (0.0030) \\
2000 & 2000 & 200 & 0.1971 (0.0028) & 0.1965 (0.0029) & 0.1934 (0.0034) \\
\bottomrule
\end{tabular}
\end{table}

Table~\ref{tab:simulation-accuracy} shows that distributed estimation achieves
accuracy close to the pooled references. At the default sample sizes, GD,
RF, and Kernel attain mean RMSEs of $0.2075$, $0.2064$, and $0.2061$,
respectively. Increasing the source sample size from $1000$ to $8000$
reduces GD RMSE from $0.2328$ to $0.1947$, while increasing the target
sample size from $250$ to $2000$ reduces it from $0.2836$ to $0.1971$.
Both source and target observations therefore improve accuracy, with smaller
gains at the larger sample sizes considered. GD nevertheless has a slightly
higher mean RMSE than RF in these comparisons. For example, at
$(N,M)=(1000,1000)$, the paired mean RMSE of GD exceeds that of RF by
$0.0036$, with a $95\%$ bootstrap confidence interval of
$[0.0003,0.0072]$ based on $2000$ dataset-level resamples. Thus the methods
have similar accuracy, with a small observed disadvantage for GD in this
setting. Because GD and RF select their parameters separately, the difference
can reflect both finite optimization and parameter selection. The next
experiment uses identical parameters to isolate optimization error.

\subsection{Optimization and the Accuracy--Communication Trade-off}
\label{Sec: simulation-communication}

At the default sample sizes, we vary $R\in\{50,200,400\}$ and
$\lambda\in\{0.001,0.01,0.1\}$, fixing $\sigma=\sigma_0$.
For each replication and parameter combination, GD and RF use identical
data and feature maps. Features are nested across $R$ using a common master
draw of frequencies and phases. We record the objective gap relative to
the exact RF solution, prediction discrepancies, and test RMSE along the
GD trajectory, through $20000$ updates. We also compare feature dimensions
at fixed communication budgets
$W\in\{10^5,4\times10^5,1.6\times10^6,6.4\times10^6\}$, allowing
$T=\lfloor W/(2CR)\rfloor$ updates.

Communication is measured in transmitted scalar coordinates. GD fitting
requires $2CRT$ coordinates, whereas transmitting the upper triangle of
each symmetric client second-moment matrix requires $CR(R+1)/2$.
These are fitting costs: initialization, validation and parameter tuning,
and optional final-model delivery are accounted for separately. The
experiment simulates the client/server interface on one computer and does
not measure network latency.

\begin{table}[htbp]
\centering
\caption{Descriptive communication comparison at $\lambda=0.01$ over
$100$ replications. $T_{2\%}$ is the first recorded checkpoint whose mean
GD test RMSE is no greater than $1.02$ times the mean RF test RMSE.
Costs are scalar coordinates; the final column is GD cost divided
by symmetric-matrix cost. This test-based summary is not a stopping rule.}
\label{tab:simulation-communication}
\small
\begin{tabular}{rcrrrr}
\toprule
$R$ & RF RMSE & $T_{2\%}$ & Matrix cost & GD cost & Cost ratio \\
\midrule
50  & 0.2053 & 500 & 5100   & 200000 & 39.22 \\
200 & 0.1891 & 200 & 80400  & 320000 & 3.98 \\
400 & 0.1859 & 200 & 320800 & 640000 & 2.00 \\
\bottomrule
\end{tabular}
\end{table}

The trajectories approach the matched RF solution as optimization proceeds.
Table~\ref{tab:simulation-communication} illustrates the resulting trade-off
at $\lambda=0.01$. Increasing $R$ improves the closed-form reference RMSE
from $0.2053$ to $0.1859$, but increases the cost of every GD round.
Consequently, a larger feature representation must be balanced against
the number of affordable updates.

Vector exchanges reduce the payload of each round from a quadratic to a
linear dependence on $R$, but this does not imply a smaller cumulative
payload. In the displayed settings, approaching the RF mean RMSE within
$2\%$ costs more than one-time symmetric-matrix aggregation. The empirical
benefit is therefore an adjustable, vector-only optimization procedure with
an explicit accuracy--communication trade-off; these results do not establish
a uniform reduction in total communication.

\subsection{Heterogeneous Client Laws and Unequal Client Sizes}
\label{Sec: simulation-heterogeneity}

We first vary client heterogeneity while keeping the pooled estimation
problem fixed. With $\pi_c=1/4$, define
\[
P_c^{(\tau)}=(1-\tau)P+\tau F_c,
\qquad \tau\in\{0,0.5,1\}.
\]
Then $\sum_c\pi_cP_c^{(\tau)}=P$ for every $\tau$, while $Q$ is unchanged.
Clients range from identically distributed at $\tau=0$ to
component-specific at $\tau=1$. We fix $\lambda=0.01$,
$\sigma=\sigma_0$, $R=200$, and $T=1000$. Kernel uses the same bandwidth
and regularization. This matched-parameter comparison isolates
heterogeneity rather than independently optimizing each method.

We additionally consider naive local KuLSIF: client $c$ estimates $q/p_c$
using its own source observations and the target sample. For this diagnostic,
the prediction at a test point drawn from client $c$ is that client's fitted
ratio. Its error against the pooled truth is aggregated using $\pi_c$.
Thus the local procedure supplies client-dependent predictions and targets
different ratios; it is not an estimator of a single pooled ratio.

\begin{table}[htbp]
\centering
\caption{Controlled heterogeneity with fixed pooled $P$ and target $Q$:
mean RMSE with MCSE in parentheses, over $100$ replications per setting.
Local denotes the client-specific diagnostic described in the text.}
\label{tab:simulation-heterogeneity}
\small
\begin{tabular}{ccccc}
\toprule
$\tau$ & GD & RF & Kernel & Local \\
\midrule
0   & 0.1915 (0.0020) & 0.1915 (0.0020) & 0.1878 (0.0019) & 0.2163 (0.0018) \\
0.5 & 0.1900 (0.0016) & 0.1900 (0.0016) & 0.1839 (0.0017) & 0.5525 (0.0032) \\
1   & 0.1894 (0.0018) & 0.1894 (0.0018) & 0.1845 (0.0018) & 1.1578 (0.0044) \\
\bottomrule
\end{tabular}
\end{table}

Table~\ref{tab:simulation-heterogeneity} shows stable distributed accuracy:
GD RMSE remains close to $0.19$ throughout and agrees with RF to the
reported precision. In contrast, the local diagnostic error increases
from $0.2163$ to $1.1578$. This contrast illustrates the importance of
estimating the pooled ratio: heterogeneity need not substantially impair
the proposed estimator, whereas client-specific ratios increasingly differ
from the required pooled importance weights. Small changes across $\tau$
are possible because the stratified sampling variance also changes.

Finally, we set $P_c=F_c$ and
$\pi=(0.7,0.1,0.1,0.1)$, retaining $(N,M)=(2000,1000)$ and the same $Q$.
We recompute the pooled truth, whose upper bound is now $5$, and use the
validation-based tuning protocol. Over $500$ replications, GD achieves
RMSE $0.4197$ (MCSE $0.0036$), compared with $0.4232$ ($0.0040$) for RF
and $0.4292$ ($0.0056$) for Kernel. Replacing the correct client weights
with uniform weights, while retaining GD's selected parameters, feature
map, and iteration count, increases RMSE to $0.8968$ ($0.0017$).
This controlled ablation demonstrates the practical importance of
sample-proportion weighting. The balanced and unequal-size designs have
different pooled source laws and true ratios, so their difference in RMSE
cannot be attributed solely to unequal client sample sizes.

\paragraph{Additional diagnostic and scope.}
We also evaluated the raw ratio of pooled Gaussian kernel density estimates,
selecting source and target bandwidths by the same validation criterion over
the nine pairs of $\{0.5,1,2\}$ times their respective fitting-data
normal-reference bandwidths in the tuned settings. Controlled-heterogeneity
settings used the normal-reference bandwidths without tuning. Of $3206$
selected KDE fits, $356$ have provably infinite conditional population
$L^2(P)$ risk under the Gaussian-tail diagnostic, despite the bounded true
ratio. Their finite-test errors are therefore descriptive and are not used
to claim general superiority over KDE. No denominator floor or clipping
was introduced. Overall, the three experiments support accurate pooled-ratio
estimation across the client configurations studied and quantify its
optimization--communication trade-off. The conclusions concern this
five-dimensional mixture family; they do not establish general robustness
or verify the theoretical convergence rates.
